\section{Simulation Study}\label{sec:simulation}

This section will present the four main simulation scenarios, summarize the outputs, and describe the core figures and tables.

% NOTE: Simulation grid uses four imbalance patterns:
% - balanced: equal group sizes
% - mild: shallow deterministic linear weights (1.2 to 0.8)
% - strong: steep deterministic linear weights (G to 1)
% - random: group proportions drawn from a flat Dirichlet(1,...,1),
%   equivalent to drawing G independent Exponential(1) values and
%   normalizing by their sum. This pattern represents unsystematic
%   random imbalance with no preferred direction or ordering, and
%   samples the full range from near-balance to extreme concentration.
%   Each replicate redraws the random proportions, so the "random"
%   grid rows reflect average behavior over the entire imbalance
%   spectrum. The Gamma(1,1) / flat Dirichlet choice is supported
%   by the fact that it is distribution-free over the simplex: every
%   partition of n across G groups receives equal prior probability.
%   Document this justification in the simulation write-up.
%
% NOTE 2026-07-16: The grid includes degenerate settings at G=40, n=100
% (min_group_size = 2).  With only 2 observations per group, the
% subgroup OLS model (intercept + x1 + x2) is nearly singular,
% producing extreme coefficient estimates and inflating drift into
% the hundreds or thousands.  These 6 rows (out of 140) skew
% pattern-level means dramatically: e.g. mild mean_drift jumps from
% ~3 (median) to ~1129 (mean) because G=40 n=100 mild produced
% drift ≈ 39303 in one random draw.  All aggregated comparisons in
% the paper should either (a) use medians or (b) explicitly exclude
% settings with min_group_size < 3 (dropping 6 rows, leaving 134).
% Context-drift follows the same pattern: median ctx-drift is
% well-behaved across patterns, while means are inflated by the same
% degenerate rows.